\section{Further Steps Towards Rank Preservation} \label{section:rank_preservation}

The ability to preserve DFT-calculated interface rankings is critical for efficient high-throughput screening of material systems. Absolute interface energies are less relevant in isolation if the relative favourability of competing configurations is misrepresented. This motivates the need for improving rank fidelity in machine-learned models such as MACE.

Future work will systematically compare three training paradigms to address limitations identified in Chapter 3 :

This baseline model is trained on broad-spectrum, bulk-focused data. While it demonstrates some transferability, its limitations are apparent in systems with interfacial complexity; especially silicon-rich or structurally diverse cases.

Here, the same model is refined using additional interface-representative configurations, tailored to improve accuracy in the interfacial regime. This targets data-domain mismatch as a key source of ranking degradation.

This strategy applies a correction to existing MLP predictions by learning residuals between MLP and DFT energies. It allows the reuse of baseline MLPs while introducing physics-informed calibration, improving ranking accuracy without fully retraining the model.

% TODO: Add equation here for ΔE = E_DFT − E_MLP

\subsection{Planned Benchmarking Criteria}

Performance will be assessed using statistical correlation coefficients and overlap in Top-N rankings between DFT and MLP predictions. These metrics will provide fine-grained diagnostics of rank preservation.

% TODO: Insert table summarising benchmarking metrics (Spearman, Kendall, Top-N) and their interpretations

Benchmarks will be applied across chemically diverse systems (C, Si, Ge, Sn), crystallographic orientations, and stacking shifts. Key figures include \texttt{results\_upper\_Si\_spearman\_dft\_dft\_vs\_dft\_mlp\_valid.png} and \texttt{results\_lower\_Ge\_kendall\_dft@dft\_vs\_dft@mlp\_invalid.png}. Reference datasets include \texttt{results.csv} and \texttt{stats.csv} (1381 interfaces; no train/test partition defined).

No explicit mechanism yet exists to identify and flag phase boundary rank inversions, though these could be extracted from \texttt{\_rank\_distance} fields. Future work may focus on visualising such transitions as indicators of model degradation.

% TODO: Add proposed figure example or sketch of phase boundary inversion analysis

\paragraph{Additional Limitation:}

Error magnitude may scale with system size, particularly in Sn-rich systems. However, no current analysis addresses this dependency. A robustness study is recommended.

% TODO: Add figure showing error vs. system size if/when analysis is complete

\paragraph{Additional Limitation:}

Missing or failed calculations are marked via \texttt{is\_broken}, but no summary or completeness tracker has been implemented. Future iterations may benefit from such a tool for data transparency.